% GENERATED FILE. Edit canonical lessons, then run npm run generate:books.
% canonical-source-hash: fnv1a64:8a65d4872efb0727
% canonical-lessons: SA-C515-svadah, SA-C515-svaduh, SA-C515-amah, SA-C515-pakvah, SA-C515-cusati

\chapter{How Does It Taste?}
\label{ch:sa-how-does-it-taste}
\hlchaptermodality{515}

\begin{chapteropening}
\textbf{By the end of this chapter:} \emph{I can say how food tastes and whether it is cooked.}

Complete the last lesson of chapter 515: I can say how food tastes and whether it is cooked.
\end{chapteropening}

\section[\texorpdfstring{\sk{स्वादः} (svādaḥ)}{स्वादः (svādaḥ)}]{\sk{स्वादः} (svādaḥ) — a flavour}

\label{lesson:SA-C515-svadah}



\begin{quote}
Before the new one: say the Sanskrit for leftovers, then the Sanskrit for a lid.
\end{quote}

\subsection*{You'll want to know: \sk{स्वादः}}
\textbf{\sk{स्वादः}} — \emph{svādaḥ} — \textquotedblleft{}a flavour\textquotedblright{}.

\begin{quote}
\sk{एतस्य} \sk{सूपस्य} \sk{स्वादः} \sk{उत्तमः} \sk{अस्ति।}
\end{quote}

\begin{quote}
\emph{etasya sūpasya svādaḥ uttamaḥ asti.}
\end{quote}

\begin{quote}
This dal has a fine flavour.
\end{quote}

A masculine noun.

\begin{cousinweb}
\textbf{\sk{स्वादः}} (\emph{svādaḥ}) comes from the root √svad, \textquotedblleft{}to be sweet, to relish\textquotedblright{}; MW gives \emph{svāda} as taste, flavour, savour (MW).
\end{cousinweb}

\begin{grammarlens}[title={What you've built}]
One more word for describing what you lived and felt.
\end{grammarlens}

\subsection*{Your turn}
\begin{itemize}
\raggedright
  \item \emph{Say it:} \emph{svādaḥ}
  \item \emph{Say it:} \emph{svādaḥ}, once more
  \item \emph{Say it:} \emph{etasya sūpasya svādaḥ uttamaḥ asti.} — This dal has a fine flavour.
  \item \emph{Recall:} say the Sanskrit for leftovers, then the Sanskrit for a lid, then say \emph{svādaḥ} again
\end{itemize}

\begin{tcolorbox}[breakable,colback=teal!4,colframe=teal!35!black,title={Before you move on}]
What does \sk{स्वादः} mean? (\textquotedblleft{}A flavour\textquotedblright{}.) Say it once more.
\end{tcolorbox}
\section[\texorpdfstring{\sk{स्वादुः} (svāduḥ)}{स्वादुः (svāduḥ)}]{\sk{स्वादुः} (svāduḥ) — tasty}

\label{lesson:SA-C515-svaduh}



\begin{quote}
Before the new one: say the Sanskrit for a lid, then the Sanskrit for a flavour.
\end{quote}

\subsection*{You'll want to know: \sk{स्वादुः}}
\textbf{\sk{स्वादुः}} — \emph{svāduḥ} — \textquotedblleft{}tasty\textquotedblright{}.

\begin{quote}
\sk{मातुः} \sk{भोजनं} \sk{स्वादु} \sk{अस्ति।}
\end{quote}

\begin{quote}
\emph{mātuḥ bhojanaṃ svādu asti.}
\end{quote}

\begin{quote}
Mother's cooking is tasty.
\end{quote}

\begin{cousinweb}
\textbf{\sk{स्वादुः}} (\emph{svāduḥ}) comes from the root √svad, \textquotedblleft{}to be sweet\textquotedblright{}; MW gives \emph{svādu} as sweet, savoury, pleasant to the taste (MW).
\end{cousinweb}

\begin{grammarlens}[title={What you've built}]
One more word for describing what you lived and felt.
\end{grammarlens}

\subsection*{Your turn}
\begin{itemize}
\raggedright
  \item \emph{Say it:} \emph{svāduḥ}
  \item \emph{Say it:} \emph{svāduḥ}, once more
  \item \emph{Say it:} \emph{mātuḥ bhojanaṃ svādu asti.} — Mother's cooking is tasty.
  \item \emph{Recall:} say the Sanskrit for a lid, then the Sanskrit for a flavour, then say \emph{svāduḥ} again
\end{itemize}

\begin{tcolorbox}[breakable,colback=teal!4,colframe=teal!35!black,title={Before you move on}]
What does \sk{स्वादुः} mean? (\textquotedblleft{}Tasty\textquotedblright{}.) Say it once more.
\end{tcolorbox}
\section[\texorpdfstring{\sk{आमः} (āmaḥ)}{आमः (āmaḥ)}]{\sk{आमः} (āmaḥ) — raw; uncooked}

\label{lesson:SA-C515-amah}



\begin{quote}
Before the new one: say the Sanskrit for a flavour, then the Sanskrit for tasty.
\end{quote}

\subsection*{You'll want to know: \sk{आमः}}
\textbf{\sk{आमः}} — \emph{āmaḥ} — \textquotedblleft{}raw; uncooked\textquotedblright{}.

\begin{quote}
\sk{एतत्} \sk{शाकम्} \sk{आमम्} \sk{अस्ति।}
\end{quote}

\begin{quote}
\emph{etat śākam āmam asti.}
\end{quote}

\begin{quote}
This vegetable is raw.
\end{quote}

\begin{cousinweb}
\textbf{\sk{आमः}} (\emph{āmaḥ}) is the adjective MW records as raw, uncooked, the opposite of \emph{pakva}, \textquotedblleft{}cooked\textquotedblright{}; MW gives no derivation, so its origin is uncertain (MW).
\end{cousinweb}

\begin{grammarlens}[title={What you've built}]
One more word for describing what you lived and felt.
\end{grammarlens}

\subsection*{Your turn}
\begin{itemize}
\raggedright
  \item \emph{Say it:} \emph{āmaḥ}
  \item \emph{Say it:} \emph{āmaḥ}, once more
  \item \emph{Say it:} \emph{etat śākam āmam asti.} — This vegetable is raw.
  \item \emph{Recall:} say the Sanskrit for a flavour, then the Sanskrit for tasty, then say \emph{āmaḥ} again
\end{itemize}

\begin{tcolorbox}[breakable,colback=teal!4,colframe=teal!35!black,title={Before you move on}]
What does \sk{आमः} mean? (\textquotedblleft{}Raw; uncooked\textquotedblright{}.) Say it once more.
\end{tcolorbox}
\section[\texorpdfstring{\sk{पक्वः} (pakvaḥ)}{पक्वः (pakvaḥ)}]{\sk{पक्वः} (pakvaḥ) — ripe; cooked}

\label{lesson:SA-C515-pakvah}



\begin{quote}
Before the new one: say the Sanskrit for tasty, then the Sanskrit for raw.
\end{quote}

\subsection*{You'll want to know: \sk{पक्वः}}
\textbf{\sk{पक्वः}} — \emph{pakvaḥ} — \textquotedblleft{}ripe; cooked\textquotedblright{}.

\begin{quote}
\sk{एतत्} \sk{आम्रं} \sk{पक्वं} \sk{न} \sk{अस्ति।}
\end{quote}

\begin{quote}
\emph{etat āmraṃ pakvaṃ na asti.}
\end{quote}

\begin{quote}
This mango is not ripe.
\end{quote}

\begin{cousinweb}
\textbf{\sk{पक्वः}} (\emph{pakvaḥ}) comes from the root √pac, \textquotedblleft{}to cook\textquotedblright{}, and MW treats it as that root's past participle: cooked, baked, and so ripe (MW).
\end{cousinweb}

\begin{grammarlens}[title={What you've built}]
One more word for describing what you lived and felt.
\end{grammarlens}

\subsection*{Your turn}
\begin{itemize}
\raggedright
  \item \emph{Say it:} \emph{pakvaḥ}
  \item \emph{Say it:} \emph{pakvaḥ}, once more
  \item \emph{Say it:} \emph{etat āmraṃ pakvaṃ na asti.} — This mango is not ripe.
  \item \emph{Recall:} say the Sanskrit for tasty, then the Sanskrit for raw, then say \emph{pakvaḥ} again
\end{itemize}

\begin{tcolorbox}[breakable,colback=teal!4,colframe=teal!35!black,title={Before you move on}]
What does \sk{पक्वः} mean? (\textquotedblleft{}Ripe; cooked\textquotedblright{}.) Say it once more.
\end{tcolorbox}
\section[\texorpdfstring{\sk{चूषति} (cūṣati)}{चूषति (cūṣati)}]{\sk{चूषति} (cūṣati) — sucks}

\label{lesson:SA-C515-cusati}



\begin{quote}
Before the new one: say the Sanskrit for raw, then the Sanskrit for ripe.
\end{quote}

\subsection*{You'll want to know: \sk{चूषति}}
\textbf{\sk{चूषति}} — \emph{cūṣati} — \textquotedblleft{}sucks\textquotedblright{}.

\begin{quote}
\sk{बालः} \sk{आम्रं} \sk{चूषति।}
\end{quote}

\begin{quote}
\emph{bālaḥ āmraṃ cūṣati.}
\end{quote}

\begin{quote}
The child sucks the mango.
\end{quote}

Past (he …): \sk{सः} \sk{आम्रं} \sk{चूषितवान्}, \emph{saḥ āmraṃ cūṣitavān}.

\begin{cousinweb}
\textbf{\sk{चूषति}} (\emph{cūṣati}) comes from the root √cūṣ, \textquotedblleft{}to suck, to suck out\textquotedblright{}; its past participle is \emph{cūṣita} (MW).
\end{cousinweb}

\begin{grammarlens}[title={What you've built}]
One more word for describing what you lived and felt.
\end{grammarlens}

\subsection*{Your turn}
\begin{itemize}
\raggedright
  \item \emph{Say it:} \emph{cūṣati}
  \item \emph{Say it:} \emph{cūṣati}, once more
  \item \emph{Say it:} \emph{bālaḥ āmraṃ cūṣati.} — The child sucks the mango.
  \item \emph{Recall:} say the Sanskrit for raw, then the Sanskrit for ripe, then say \emph{cūṣati} again
\end{itemize}

\begin{tcolorbox}[breakable,colback=teal!4,colframe=teal!35!black,title={Before you move on}]
What does \sk{चूषति} mean? (\textquotedblleft{}Sucks\textquotedblright{}.) Say it once more.
\end{tcolorbox}
